\section{Referencial teórico}\label{sec:referencial}

\subsection{Estudo bibliométrico}\label{sec:biblio}

A estimação dos efeitos das transferências intergovernamentais sobre o gasto e a arrecadação dos governos subnacionais está no centro da literatura de federalismo fiscal. As transferências financiam boa parte da provisão local de serviços em países federais e unitários descentralizados, e a forma como os governos receptores as utilizam --- gastando-as, reduzindo tributos ou substituindo esforço próprio --- é decisiva para o desenho dos sistemas de repartição. Duas regularidades organizam esse debate: o efeito \textit{flypaper}, segundo o qual as transferências elevam o gasto muito mais do que um aumento equivalente da renda local, e a ilusão fiscal, segundo a qual o financiamento por terceiros distorce a percepção do custo dos serviços públicos.

A Tabela~\ref{tab:baker}, construída no mesmo formato da Tabela 1 de \citeonline{baker2022}, mostra que, de 2000 a 2025, foram publicados \tUmTot{} artigos sobre transferências intergovernamentais nos quinze periódicos selecionados: \tUmFin{} em periódicos de finanças públicas e escolha pública, \tUmFed{} em periódicos de federalismo, orçamento e governo local e \tUmBr{} em periódicos brasileiros. Entre eles, \tUmFly{} (\tUmPct\%) tratam do efeito \textit{flypaper} ou da ilusão fiscal, e \tUmDezTot{} desses \tUmFly{} (\tUmDezPct\%) foram publicados a partir de 2010. No corpus completo do período, com \corpTot{} artigos em todos os periódicos, a participação é semelhante (\corpFly{} artigos, ou \corpPct\%), e \corpDezPct\% deles são posteriores a 2010. A participação dos dois temas é alta em todos os grupos, inclusive nos periódicos brasileiros (\tUmBrPct\%), o que mostra que o efeito \textit{flypaper} e a ilusão fiscal não são tópicos marginais, mas um eixo central da pesquisa sobre transferências.

\begin{table}[htbp]
\centering
\caption{Artigos sobre transferências intergovernamentais e sobre efeito \textit{flypaper} ou ilusão fiscal, por periódico: 2000--2025}
\label{tab:baker}
\small
\begin{tabular}{lccc}
\toprule
 & (1) & (2) & (3)\\
 & Transferências & \textit{Flypaper}/ilusão fiscal & (2)/(1) (\%)\\
\midrule
\textit{Public Finance Review} & 26 & 13 & 50,0\%\\
\textit{International Tax and Public Finance} & 24 & 5 & 20,8\%\\
\textit{Public Choice} & 22 & 14 & 63,6\%\\
\textit{Journal of Public Economics} & 12 & 3 & 25,0\%\\
\textit{National Tax Journal} & 7 & 4 & 57,1\%\\
\textbf{Finanças públicas e escolha pública} & \textbf{91} & \textbf{39} & \textbf{42,9\%}\\
\midrule
\textit{Public Budgeting \& Finance} & 17 & 10 & 58,8\%\\
\textit{Publius: The Journal of Federalism} & 15 & 3 & 20,0\%\\
\textit{Local Government Studies} & 12 & 7 & 58,3\%\\
\textit{Public Finance and Management} & 13 & 4 & 30,8\%\\
\textit{State and Local Government Review} & 7 & 1 & 14,3\%\\
\textbf{Federalismo, orçamento e governo local} & \textbf{64} & \textbf{25} & \textbf{39,1\%}\\
\midrule
\textit{Economia Aplicada} & 7 & 3 & 42,9\%\\
\textit{Revista de Administração Pública} & 6 & 0 & 0,0\%\\
\textit{Revista Econômica do Nordeste} & 6 & 1 & 16,7\%\\
\textit{Administração Pública e Gestão Social} & 5 & 4 & 80,0\%\\
\textit{Estudos Econômicos} & 4 & 2 & 50,0\%\\
\textbf{Periódicos brasileiros} & \textbf{28} & \textbf{10} & \textbf{35,7\%}\\
\midrule
\textbf{Total dos 15 periódicos} & \textbf{183} & \textbf{74} & \textbf{40,4\%}\\
Corpus completo (todos os periódicos) & 1.712 & 640 & 37,4\%\\
\bottomrule
\end{tabular}
\par\smallskip
\begin{minipage}{0.97\textwidth}\scriptsize
Nota: a Tabela~\ref{tab:baker} segue o formato da Tabela 1 de \citeonline{baker2022}. Resume o número de artigos publicados entre 2000 e 2025 em cinco periódicos de finanças públicas e escolha pública, cinco periódicos de federalismo, orçamento e governo local (os mais produtivos do corpus em cada grupo) e cinco periódicos editados no Brasil. O corpus foi obtido no OpenAlex com busca em título e resumo pelos termos ``intergovernmental transfers'', ``intergovernmental grants'', ``flypaper effect'' e ``fiscal illusion'' e pelos equivalentes em português e espanhol (``transferências intergovernamentais'', ``transferências constitucionais'', ``efeito flypaper'', ``ilusão fiscal'', ``transferencias intergubernamentales'', ``ilusión fiscal''), mantendo apenas artigos com autoria identificada e fonte primária classificada como periódico. A coluna (1) informa o total de artigos do corpus em cada periódico; a coluna (2), quantos deles tratam do efeito \textit{flypaper} ou da ilusão fiscal (termo presente no título, no resumo ou nas palavras-chave); a coluna (3), a participação de (2) em (1). Contagens produzidas com o pacote \textit{bibliometrix} \cite{aria2017} (\texttt{06\_codigos\_e\_planilhas/05\_bibliometria.R}).
\end{minipage}
\end{table}


A difusão desses estudos reflete a percepção de que o efeito \textit{flypaper} é uma das anomalias mais robustas da economia do setor público \cite{hines1995}. Ainda assim, avanços recentes da econometria aplicada mostram que estratégias de identificação muito usadas podem produzir estimativas enganosas. \citeonline{goodmanbacon2021} demonstra que o estimador de efeitos fixos bidirecionais com adoção escalonada é uma média ponderada de comparações entre pares de grupos, algumas delas usando unidades já tratadas como controle; \citeonline{dechaisemartin2020} mostram que os pesos dessas comparações podem ser negativos quando os efeitos são heterogêneos; \citeonline{callaway2021} propõem estimadores por coorte que contornam o problema; e \citeonline{baker2022} mostram que esses vieses podem inverter conclusões publicadas. No corpus analisado, \nDid{} artigos mencionam diferenças-em-diferenças no título, no resumo ou nas palavras-chave, \nDidQuinze{} deles a partir de 2015, o que indica que a literatura sobre transferências começa a adotar esses desenhos. Como as transferências municipais brasileiras não constituem um tratamento discreto adotado em datas distintas, este trabalho não usa diferenças-em-diferenças: estima regressões com efeito fixo de município para Manaus e Parintins e interpreta os coeficientes como associações condicionais (Seção~\ref{sec:metodo}).

O corpus foi construído no OpenAlex\footnote{OPENALEX. API de consulta bibliográfica (\textit{works}). Disponível em: \url{https://api.openalex.org}. Acesso em: 21 out. 2026.} e processado com o \textit{bibliometrix} \cite{aria2017}, por meio do conversor para a API do OpenAlex e das funções \texttt{biblioAnalysis}, \texttt{lotka}, \texttt{biblioNetwork} e \texttt{thematicMap}; o mesmo arquivo pode ser aberto na interface \texttt{biblioshiny()}. As duas consultas (em inglês e em português/espanhol) retornaram \flEn{} e \flPt{} registros, ou \flUnicos{} registros únicos; foram excluídos \flSemAut{} sem autoria, \flNaoPer{} publicados fora de periódicos e \flDoisMilVinteSeis{} de 2026, ano incompleto, restando \flFinal{} artigos (\bibAnoMin--\bibAnoMax). A produção cresce à taxa média de \bibCresc\% ao ano (Figura~\ref{fig:biblio}), \bibPctDesdeVinte\% dos artigos são posteriores a 2019, e \nBrasil{} artigos tratam do Brasil, dos quais \nBrasilFly{} abordam o efeito \textit{flypaper} ou a ilusão fiscal. A distribuição da produtividade dos autores é compatível com a lei de Lotka ($\beta=\lotkaBeta$, $R^2=\lotkaR$, $p=\lotkaP$ no teste de Kolmogorov-Smirnov), e a coautoria internacional é baixa (\bibIntl\% dos artigos). As informações gerais do corpus, a rede de coocorrência de palavras-chave e o mapa temático estão no Apêndice A.

\begin{figure}[htbp]
\centering
\caption{Produção anual de artigos sobre transferências intergovernamentais, efeito \textit{flypaper} e ilusão fiscal, 1970--\bibAnoMax}\label{fig:biblio}
\includegraphics[width=0.88\textwidth]{fig5_bibliometria_producao_anual.png}
\par\smallskip
\begin{minipage}{0.88\textwidth}\scriptsize
Fonte: elaboração própria com \textit{bibliometrix} e OpenAlex. As linhas indicam artigos cujo título, resumo ou palavras-chave mencionam o termo; artigos antigos frequentemente não têm resumo indexado, o que subestima as linhas nos primeiros anos.
\end{minipage}
\end{figure}

\subsection{Federalismo fiscal e o papel das transferências}

A teoria clássica parte da ideia de que a provisão descentralizada permite ajustar a oferta de bens públicos às preferências locais, seja pela mobilidade dos cidadãos entre jurisdições \cite{tiebout1956}, seja pela vantagem informacional dos governos próximos à população \cite{oates1999}. Nesse arcabouço, as transferências cumprem funções bem definidas: corrigir externalidades interjurisdicionais, equalizar capacidades fiscais e cobrir o desequilíbrio vertical decorrente da centralização das bases tributárias mais eficientes \cite{bird2002,bahl2008}. A segunda geração do federalismo fiscal, sintetizada por \citeonline{weingast2009} e \citeonline{oates2008}, incorporou os incentivos políticos e fiscais: governos subnacionais que dependem de transferências têm menos estímulo para promover a base econômica local e mais incentivo para transferir custos ao governo central. A evidência comparada de \citeonline{rodden2002} mostra que déficits persistentes surgem quando os entes subnacionais dependem de transferências gerais e têm liberdade para se endividar, e \citeonline{rodden2003} conclui que a descentralização só contém o crescimento do governo quando é financiada por tributos próprios, e não por repasses. Na América Latina, \citeonline{stein1999} associa o desequilíbrio vertical ao problema do fundo comum, que amplia o tamanho do governo, e os modelos de economia política do orçamento explicam por que a dispersão do custo e a concentração do benefício geram viés deficitário \cite{alesina1995,vonhagen1995}.

A literatura também mostra que a alocação das transferências não é neutra. Governos centrais usam repasses discricionários de forma tática para favorecer eleitores decisivos \cite{johansson2003} e prefeitos alinhados politicamente, como documentado para a Espanha \cite{soleolle2008,simon2013}, para o Brasil \cite{brollo2012} e em modelos de sinalização com evidência italiana \cite{bracco2015}. Mesmo as transferências constitucionais, distribuídas por fórmula, produzem efeitos que vão além do volume de gasto: no Brasil, os saltos do FPM em faixas populacionais elevaram a corrupção observada e reduziram a escolaridade média dos candidatos a prefeito \cite{brollo2013}, ao mesmo tempo em que aumentaram o gasto, a escolarização e a alfabetização e reduziram a pobreza, sem deslocar a receita própria \cite{litschig2013}, com ganhos que persistem no longo prazo \cite{carvalhofilho2022}. Auditorias federais indicam ainda que a responsabilização eleitoral reduz o desvio de recursos transferidos \cite{ferraz2011}. Esses resultados sugerem que o efeito das transferências depende do ambiente institucional em que são aplicadas, ponto retomado por \citeonline{lago2024} em ampla revisão recente.

\subsection{O efeito \textit{flypaper}: teoria, evidência e explicações}

O ponto de partida é o teorema de equivalência de \citeonline{bradford1971}: se o governo local apenas agrega as preferências do eleitor mediano, uma transferência incondicional equivale a um aumento da renda dos residentes e deveria ser gasta na mesma proporção que a renda privada, isto é, alguns centavos por real. A evidência empírica contrariou sistematicamente essa previsão. Na síntese de \citeonline{hines1995}, enquanto um real adicional de renda eleva o gasto estadual americano em 5 a 10 centavos, um real de transferência eleva-o em valores que chegam a um real --- daí a imagem de que o dinheiro ``gruda onde cai''. Estudos para municípios italianos \cite{gennari2014}, mexicanos \cite{cardenas2011}, quebequenses \cite{mehiriz2014}, alemães \cite{baskaran2015} e do estado de Wisconsin \cite{deller2005} reproduzem o fenômeno, e ele aparece também no gasto regional com saúde na Itália, acompanhado de ajuste lento \cite{levaggi2003}; no caso das transferências educacionais, a resposta depende do tipo de repasse e das regras de contrapartida \cite{fisher2000}. Mesmo desenhos quase-experimentais confirmam o efeito: a regra descontínua de transferências suecas elevou o gasto sem reduzir impostos \cite{dahlberg2008}, e os estados americanos gastaram mais de um dólar por dólar recebido do programa rodoviário de estímulo de 2009 \cite{leduc2017}. A evidência não é unânime: \citeonline{knight2002} mostra que, ao tratar a endogeneidade das transferências federais com instrumentos políticos, o efeito desaparece e há deslocamento significativo do gasto estadual; \citeonline{zampelli1986} encontra pouca evidência do efeito quando se considera a fungibilidade de transferências condicionais; e \citeonline{gamkhar1996} não encontram assimetria entre aumentos e cortes, ao passo que \citeonline{heyndels2001} identifica assimetria em municípios flamengos, com gasto que sobe quando as transferências crescem e impostos que sobem quando elas caem.

As explicações se dividem em quatro famílias. A primeira atribui o efeito a problemas de especificação e de estimação: \citeonline{fisher1982} e \citeonline{hamilton1983} questionam a comparabilidade entre renda e transferência nos modelos usuais, \citeonline{moffitt1984} mostra que restrições orçamentárias não lineares de programas com contrapartida distorcem as estimativas, \citeonline{megdal1987} demonstra, por simulação, que o MQO superestima a propensão a gastar quando a transferência é construída endogenamente, e \citeonline{becker1996} argumenta que parte do que se chama ilusão fiscal é artefato de especificação. A segunda enfatiza a assimetria de informação e o poder de agenda dos burocratas, que conhecem o tamanho da transferência e a utilizam para expandir o orçamento \cite{romer1980,filimon1982}; \citeonline{strumpf1998} deriva desse arcabouço um índice preditivo da magnitude do efeito, e \citeonline{wyckoff1991} confronta empiricamente as explicações concorrentes. A terceira explica o efeito pelo custo marginal dos recursos públicos: como os tributos locais geram peso morto, um real transferido custa menos ao eleitor do que um real arrecadado, o que torna racional gastar mais da transferência \cite{hamilton1986,dahlby2011}; o custo elevado da arrecadação local amplia a resposta às transferências em municípios peruanos \cite{aragon2013}. A quarta família, comportamental, associa o efeito à contabilidade mental --- recursos ``do grupo'' são tratados como distintos dos recursos individuais, como mostram experimentos de laboratório \cite{becker2020} --- e à incerteza: quando as transferências são menos voláteis ou menos correlacionadas com a renda privada, a poupança precaucional responde de modo diferente aos dois choques \cite{vegh2015}. Há ainda interpretações institucionais, segundo as quais as instituições políticas moderam a resposta às transferências e ajudam a explicar a evidência mista \cite{kjaergaard2015}, sendo o efeito menor onde a competição política e a eficiência arrecadatória são maiores \cite{rios2021}, e a constatação de que as transferências de equalização se capitalizam no preço dos imóveis, o que altera a leitura do efeito \cite{allers2016}. Revisões da literatura mantêm o diagnóstico de que o efeito existe, mas sua magnitude depende do desenho das transferências e da estratégia de identificação \cite{bailey1998,lago2024}.

\subsection{Ilusão fiscal}

A hipótese de ilusão fiscal sustenta que os contribuintes subestimam o custo dos serviços públicos quando a estrutura de receitas é complexa, pouco visível ou financiada por terceiros \cite{wagner1976}. O estudo pioneiro de \citeonline{pommerehne1978} para cidades suíças encontrou gastos maiores onde a carga tributária era menos perceptível, resultado que estimulou testes para a complexidade tributária \cite{misiolek1988,heyndels1995}, para municípios flamengos \cite{heyndels1994}, para o Reino Unido \cite{gemmell1999} e para a dívida pública local \cite{banzhaf2013}. \citeonline{dollery1996} organizam a literatura em cinco hipóteses --- complexidade da receita, elasticidade da receita, efeito \textit{flypaper}, ilusão do inquilino e ilusão da dívida ---, das quais a terceira é a que interessa diretamente a este trabalho. Nela, as transferências reduzem o preço percebido do bem público local \cite{logan1986,turnbull1992,turnbull1998}, e o contribuinte que não distingue os recursos recebidos de outras esferas tende a apoiar orçamentos maiores; \citeonline{logan1986} mostra que, simetricamente, o governo que concede a transferência parece mais caro, efeito confirmado com dados de pesquisa por \citeonline{hewitt1986}. A ilusão é maior quando a transferência vem de uma esfera mais distante \cite{grossman1990}, e \citeonline{mitias2001} identificam que o gasto dos condados americanos responde à ilusão de transferência, e não à ilusão tributária.

A evidência é heterogênea e alimenta um debate sobre a própria existência da ilusão. Experimentos de laboratório indicam que tributos menos visíveis são subestimados pelos participantes e distorcem suas escolhas de gasto \cite{sausgruber2005}, e \citeonline{baekgaard2016} mostram que ela decorre menos da falta de informação do que da falta de atenção. Índices construídos por modelos de múltiplos indicadores mostram grande variação entre países \cite{dellanno2012} e na União Europeia \cite{dellanno2014}. Para governos locais, o fenômeno aparece em municípios australianos \cite{worthington1999} e sul-africanos, em que a forte dependência de transferências é vista como fonte de ilusão \cite{amusa2008}. No Brasil, há evidência de ilusão fiscal associada a transferências e déficits nos estados \cite{prado2020} e na demanda por gasto federal \cite{silva2015}; nos municípios, a captura de transferências por grupos de interesse cresce com a ilusão fiscal e com o sobrefinanciamento de alguns entes \cite{mendes2005}, e a natureza das transferências relaciona-se com o cumprimento das metas fiscais \cite{araujo2020}.

\subsection{Evidências para os municípios brasileiros}

O estudo de \citeonline{cossio2001} é uma das primeiras evidências do efeito \textit{flypaper} no Brasil: as transferências constitucionais de natureza incondicional têm efeito mais expansivo sobre as despesas municipais do que aumentos equivalentes na renda, com transbordamentos espaciais entre municípios vizinhos. \citeonline{mattos2011} reinterpretam o fenômeno como perda de eficiência arrecadatória, e não apenas como expansão do gasto, ao mostrar que municípios que recebem mais transferências exploram menos sua base tributária; \citeonline{sakurai2013} encontra respostas assimétricas a aumentos e reduções de transferências, com recomposição entre gasto corrente e investimento; e \citeonline{nojosa2018} mostram que o efeito é mais intenso onde o prefeito tem maior força política na câmara. \citeonline{ferreira2019} confirmam o efeito em municípios brasileiros com instrumentos de alinhamento político, intensificado quando prefeito, governador e presidente são do mesmo partido, e \citeonline{ferreira2022} constroem um índice do efeito para estados e municípios entre 2000 e 2018, documentando sua relevância generalizada. A tese de Ferreira (2021)\footnote{FERREIRA, L. O. G. \textit{Constitutional tax transfer in Brazilian municipalities: the fiscal illusion and the flypaper effect}. 2021. Tese (Doutorado) -- Universidade de Brasília, Brasília, 2021.} articula esses resultados à ilusão fiscal produzida pelas transferências constitucionais.

Outros estudos qualificam o mecanismo. As transferências em bloco ampliam a dispersão do gasto educacional entre municípios, ao contrário das condicionais \cite{mattos2018}; a regra de gasto mínimo em educação elevou o gasto apenas nos municípios que estavam abaixo do piso \cite{cruz2020}; parte do efeito \textit{flypaper} decorre de transbordamentos das transferências recebidas por municípios vizinhos \cite{castro2021}; as faixas do FPM afetam a composição do gasto e os incentivos fiscais locais \cite{alves2024}; e a descentralização altera a composição da despesa de forma heterogênea ao longo da distribuição \cite{sousa2021}. Receitas de royalties, que funcionam como transferências não vinculadas, reduziram o esforço tributário de municípios produtores de petróleo \cite{postali2015}, o que reforça a hipótese de substituição. Na saúde, o financiamento federal cresceu sobretudo por transferências, que superaram 70\% do gasto federal no setor na década de 2000 \cite{machado2014}. Do ponto de vista do gasto agregado, \citeonline{mello2002} relaciona as fontes de financiamento à composição do gasto local e ao crescimento.

A disciplina fiscal municipal brasileira se apoia na Lei de Responsabilidade Fiscal\footnote{BRASIL. Lei Complementar n.º 101, de 4 de maio de 2000 (Lei de Responsabilidade Fiscal). Disponível em: \url{https://www.planalto.gov.br/ccivil_03/leis/lcp/lcp101.htm}.}, cujos efeitos sobre o endividamento municipal foram documentados para os municípios com mais de 100 mil habitantes \cite{macedo2009} e para o Piauí \cite{linhares2013}. A lei se insere na agenda de administração gerencial que buscou deslocar o foco do controle de processos para o de resultados \cite{bresserpereira1996}\footnote{BRASIL. Presidência da República. Câmara da Reforma do Estado. \textit{Plano Diretor da Reforma do Aparelho do Estado}. Brasília, 1995.}. A literatura bibliométrica nacional sobre transferências, por sua vez, concentra-se nas transferências da União para os municípios e na aplicação dos recursos em saúde e educação \cite{gomes2012}.

\subsection{A experiência espanhola}

A Espanha construiu, a partir de 1978, um Estado autonômico em que as Comunidades Autônomas assumiram saúde, educação e serviços sociais\footnote{ESPAÑA. \textit{Constitución Española de 1978}. Disponível em: \url{https://www.boe.es/eli/es/c/1978/12/27/(1)/con}.}. Até as reformas de 1996 e 2001, as regiões de regime comum eram financiadas majoritariamente por transferências; desde então, passaram a receber parcelas cedidas do IRPF, do IVA e dos impostos especiais, com capacidade normativa limitada \cite{toboso2010}. A reforma de 2009 reorganizou os fundos de garantia e de suficiência \cite{zabalza2011}\footnote{ESPAÑA. Ley 22/2009, de 18 de diciembre, que regula o sistema de financiação das Comunidades Autónomas de regime comum. Disponível em: \url{https://www.boe.es/eli/es/l/2009/12/18/22/con}.}, e a maior autonomia tributária endureceu a restrição orçamentária regional, reduzindo o incentivo a acumular déficits \cite{arespa2022}. Esse resultado converge com a evidência para a OCDE de que a descentralização ``parcial'', financiada por transferências, compromete a disciplina fiscal subnacional, ao passo que a autonomia tributária a fortalece \cite{asatryan2015}. A descentralização também se associa ao tamanho dos governos regionais espanhóis \cite{cantarero2012}, e a distribuição das transferências discricionárias responde a estratégias eleitorais \cite{simon2013}.

No nível local, os municípios espanhóis combinam tributos próprios (IBI, IAE, imposto sobre veículos e sobre construções) com a participação nos tributos do Estado\footnote{ESPAÑA. Real Decreto Legislativo 2/2004, de 5 de marzo, Texto Refundido de la Ley Reguladora de las Haciendas Locales. Disponível em: \url{https://www.boe.es/eli/es/rdlg/2004/03/05/2/con}.}. A fórmula de repartição pondera a população com pesos crescentes, refletindo custos maiores dos municípios grandes \cite{soleolle2005}; a expansão urbana dispersa elevou receitas vinculadas à construção e despesas de infraestrutura \cite{hortasrico2014}; e a situação financeira municipal se explica tanto por fatores políticos quanto socioeconômicos \cite{bastida2009}. Para 2.451 municípios, \citeonline{rios2021} mostram que o efeito \textit{flypaper} é menor onde a eficiência arrecadatória e a força política do governo são maiores, enquanto a fragmentação partidária reduz a estabilidade dos governos locais \cite{carozzi2022}. Desde 2012, a Lei Orgânica de Estabilidade Orçamentária impõe a todas as administrações limites de déficit, dívida e uma regra de gasto\footnote{ESPAÑA. Ley Orgánica 2/2012, de 27 de abril, de Estabilidad Presupuestaria y Sostenibilidad Financiera. Disponível em: \url{https://www.boe.es/eli/es/lo/2012/04/27/2/con}.}, articulada ao Pacto de Estabilidade e Crescimento e ao novo arcabouço de governança econômica europeu\footnote{UNIÃO EUROPEIA. Regulamento (UE) 2024/1263 relativo à coordenação eficaz das políticas económicas e à supervisão orçamental multilateral. \textit{Jornal Oficial da União Europeia}, L 2024/1263.}. O ajuste fiscal das administrações locais depende, em boa medida, de como as transferências reagem a choques de receita e despesa, como mostra \citeonline{buettner2006} para municípios americanos, e a experiência da África Subsaariana sugere que, onde a capacidade fiscal local é limitada, a relação entre transferências e arrecadação própria não segue necessariamente o padrão de deslocamento observado em países de renda alta \cite{masaki2018}.

\subsection{Síntese, lacunas e hipóteses}

A literatura converge em três pontos. O gasto subnacional tende a responder mais às transferências do que à renda local; essa resposta é mais forte quando a base própria é estreita e cara de explorar; e a dependência de transferências enfraquece a disciplina fiscal na ausência de regras ou de autonomia tributária. Divergem, porém, as interpretações sobre se o efeito reflete ilusão dos eleitores, poder de agenda dos burocratas, custo dos recursos ou problemas de identificação. Há duas lacunas a que este trabalho procura responder: os estudos brasileiros raramente analisam municípios da Região Norte, cujas condições de arrecadação e de custo de provisão são muito particulares, e as comparações internacionais pouco exploram dados por subsetor das contas nacionais. Com base nessa discussão, são formuladas as seguintes hipóteses verificáveis:

\begin{description}
\item[H1 (dependência)] Parintins é mais dependente de transferências do que Manaus, e ambos são mais dependentes do que as Comunidades Autônomas e as entidades locais espanholas.
\item[H2 (\textit{flypaper})] Em Manaus e Parintins, a propensão marginal a gastar as transferências supera a propensão a gastar a renda local ($\beta_1>\beta_2$).
\item[H3 (ilusão fiscal e esforço tributário)] Transferências maiores associam-se a menor arrecadação tributária própria per capita ($\gamma_1<0$).
\item[H4 (disciplina fiscal)] Os níveis subnacionais espanhóis, com maior autonomia tributária e regras de estabilidade, apresentam resultados mais próximos do equilíbrio do que a administração central.
\end{description}
